\documentclass[12pt]{article}

\usepackage[top=1in, bottom=1in, left=.75in, right=.75in]{geometry}
\usepackage{amsmath,fancyhdr,verbatim}
\usepackage{graphicx}
\usepackage{txfonts}
\usepackage{multicol,coordsys}
\usepackage[scaled=0.86]{helvet}
\renewcommand{\emph}[1]{\textsf{\textbf{#1}}}
\usepackage{anyfontsize}
\usepackage[shortlabels]{enumitem}
% \usepackage{times}
% \usepackage[lf]{MinionPro}
\usepackage{tikz,pgfplots,mathrsfs}
%\def\degC{{}^\circ{\rm C}}
\def\ra{\rightarrow}
\usetikzlibrary{calc, backgrounds}
\pgfplotsset{compat = newest}
\newcommand{\blank}[1]{\rule{#1}{0.75pt}}

\pgfplotsset{my style/.append style={axis x line=middle, axis y line=
middle, xlabel={$x$}, ylabel={$y$},axis equal}}

%yticklabels={,,} , xticklabels={,,}

% \setmainfont{Times}
% \def\sansfont{Lucida Grande Bold}
\parindent 0pt
\parskip 4pt
\pagestyle{fancy}
\fancyfoot[C]{\emph{\thepage}}
\fancyhead[L]{\ifnum \value{page} > 1\relax\emph{Math 252: Final Exam}\fi}
\fancyhead[R]{\ifnum \value{page} > 1\relax\emph{Spring 2024}\fi}
\headheight 15pt
\renewcommand{\headrulewidth}{0pt}
\renewcommand{\footrulewidth}{0pt}
\let\ds\displaystyle
\def\continued{{\emph {Continued....}}}
\def\continuing{{\emph {Problem \arabic{probcount} continued....}}\par\vskip 4pt}

\newcommand{\prob}[1]{\bigskip\noindent\textbf{#1.} }
\newcommand{\pts}[1]{{\small (\textsl{#1})}}

\newcommand{\probpts}[2]{\prob{#1} \pts{#2 pts} \quad}
\newcommand{\ppartpts}[2]{\textbf{(#1)} \pts{#2 pts} \quad}
\newcommand{\epartpts}[2]{\medskip\noindent \textbf{(#1)} \pts{#2 pts} \quad}


\begin{document}
{\emph{\fontsize{26}{28}\selectfont Math F252\hfill
{\fontsize{32}{36}\selectfont Final}
\hfill Spring 2026}}

\vskip 1cm
\strut\vtop{\halign{\emph#\hskip 0.5em\hfil&#\hbox to 4in{\hrulefill}\cr
\emph{\fontsize{18}{22}\selectfont Name:}&\cr
\noalign{\vskip 10pt}}}

\vspace{0.3in}
{\fontsize{18}{22}\selectfont\emph{Rules:}}

You have \textbf{2 hours} to complete this midterm.  

Partial credit will be awarded, but you must show your work.

Calculators and books are not allowed.  You may have \textbf{one whole sheet} of letter paper, 2-sided, with notes.

Place a box around your \fbox{FINAL ANSWER} to each question.%, or use the box provided.

Turn off anything that might go beep during the exam.  Good luck!

\vfill
\def\emptybox{\hbox to 2em{\vrule height 16pt depth 8pt width 0pt\hfil}}
\def\tline{\noalign{\hrule}}
\centerline{\vbox{\offinterlineskip
{
\bf\sf\fontsize{18pt}{22pt}\selectfont
\hrule
\halign{
\vrule#&\strut\quad\hfil#\hfil\quad&\vrule#&\quad\hfil#\hfil\quad
&\vrule#&\quad\hfil#\hfil\quad&\vrule#\cr
height 3pt&\omit&&\omit&&\omit&\cr
&Problem&&Possible&&Score&\cr\tline
height 3pt&\omit&&\omit&&\omit&\cr
& 1&&  6&&\emptybox&\cr\tline
& 2&&  6&&\emptybox&\cr\tline
& 3&& 13&&\emptybox&\cr\tline
& 4&& 12&&\emptybox&\cr\tline
& 5&&  6&&\emptybox&\cr\tline
& 6&&  6&&\emptybox&\cr\tline
& 7&& 15&&\emptybox&\cr\tline
& 8&& 12&&\emptybox&\cr\tline
& 9&&  9&&\emptybox&\cr\tline
&10&& 11&&\emptybox&\cr\tline
&11&& 14&&\emptybox&\cr\tline
&12&&  9&&\emptybox&\cr\tline
&13&&  6&&\emptybox&\cr\tline
&\textsl{Extra Credit A}&& \textsl{2}&&\emptybox&\cr\tline
&\textsl{Extra Credit B}&& \textsl{2}&&\emptybox&\cr\tline
&Total&& 125&&\emptybox&\cr
}\hrule}}}

\newpage
\probpts{1}{6}  Sketch the area under the curve $f(x)=\ln x$ between $x=1$ and $x=e$.  Find this area by evaluating an integral. \\

\begin{tikzpicture}
\draw[ultra thick, ->] (-1,0) -- (5,0);
\draw[ultra thick, ->] (0,-1.5) -- (0,3.0);
\end{tikzpicture}
\vfill

\probpts{2}{6} Evaluate (compute the value of) the series $\displaystyle \sum_{n=0}^\infty \frac{(-1)^n}{3^{n+1}}$.
\vfill


\newpage
\prob{3} \ppartpts{a}{3} On the axes below, sketch the region $R$ in the first quadrant bounded by $y=\cos(x)$, $y=0$, and $x=0$. \\

\begin{tikzpicture}
\draw[ultra thick, ->] (-1,0) -- (5,0);
\draw[ultra thick, ->] (0,-1.5) -- (0,3.0);
\end{tikzpicture}

\epartpts{b}{5} Use an integral to find the volume  of the solid obtained by rotating $R$ about the $x$-axis.
\vfill

\epartpts{c}{5} Set up, but do not evaluate, an integral for the surface area of the solid in part \textbf{(b)}.
\vfill


\newpage
\prob{4}  Evaluate the indefinite integrals below.

\epartpts{a}{6} $\displaystyle{\int \cos^5(\theta) \, d\theta}$
\vfill

\epartpts{b}{6} $\displaystyle{\int \frac{1}{(4+x^2)^{3/2}} \, dx}$
\vfill


\newpage
\probpts{5}{6}  A spring with a relaxed length of 2 meters requires 3 Newtons force to stretch to a length of 2.1 meters.  How much work would it take to stretch the spring from 2 meters to 2.4 meters?
\vfill

\probpts{6}{6}  Find the Maclaurin series of\quad $\ds \sinh x = \frac{e^x - e^{-x}}{2}$.  Express your answer as \textbf{one series}.
\vfill


\newpage
\prob{7} \pts{5 pts each} \quad Do the following series converge or diverge? Show your work; \textbf{state the name} of any test you use.

\renewcommand{\labelenumi}{\textbf{(\alph{enumi})}}
	\begin{enumerate}[leftmargin=7mm]
	\item \, $\displaystyle{\sum_{n=0}^\infty \frac{(-1)^n}{\sqrt{2n+1}}}$
	\vfill
	\item \, $\displaystyle{\sum_{n=1}^\infty \frac{4^{n+1}}{n^n}}$
	\vfill
	\item \, $\displaystyle{\sum_{n=2}^\infty \frac{1}{n^2 \ln n}}$
	\vfill
	\end{enumerate}


\newpage
\prob{8} \pts{6 pts each} \quad For each power series, determine the \emph{interval} of convergence.
	\begin{enumerate}[leftmargin=7mm]
	\item \,$\displaystyle \sum_{n=1}^\infty \frac{(x+2)^n}{n}$
	\vfill
	\item \,$\displaystyle \sum_{n=0}^\infty \frac{x^{2n}}{8(n+1)!}$
	\vspace{3.0in}
	\end{enumerate}


\newpage
\prob{9} Let \,\, $\ds f(x)=\frac{1}{1+x}$. 
	\begin{enumerate}[leftmargin=7mm]
	\item \pts{4 pts} \, Find a formula for $f^{(n)}(x),$ the $n$th derivative of $f(x).$
	\vspace{3.0in}
	\item \pts{5 pts} \, Find the Taylor series for $f(x)$ centered at $a=1$. Your answer should be reasonably simplified.
	\vfill
	\end{enumerate}


\newpage
\prob{10} Consider the curve defined by the parametric equations \hspace{.2in} {\large $x=\ln(t)$, \hspace{.1in} $y=(t-1)^2$}.
	\begin{enumerate}[leftmargin=7mm]
	\item \pts{4 pts} \,  Find the equation of the tangent line at $t=2$.
	\vfill
	\item \pts{3 pts} \,  For $t>0$, determine the points on the curve---give the $x,y$ coordinates of these points---at which the tangent line is horizontal.
	\vfill
	\item \pts{4 pts} \,  Set up but do not evaluate an integral for the length of the curve from $t=1$ to $t=2.$
	\vspace{1.6in}
	\end{enumerate}  


\newpage
\prob{11}  \ppartpts{a}{4}  Find the Maclaurin series for {\large $\ds e^{-x^2}$}.  (\emph{Hint.}  Start from a known Maclaurin series.) \\
\vfill

\epartpts{b}{5}  Compute an infinite series for the integral \,\,$\ds \int_0^3 e^{-x^2}\,dx$.
\vfill

\epartpts{c}{5}  Show that the infinite series in part \textbf{(b)} converges.  \textbf{State the name} of any test you use.
\vfill


\newpage
\prob{12}  \ppartpts{a}{4} Convert the polar equation $r=\cos \theta$ to rectangular form.
\vfill

\epartpts{b}{5} Sketch the polar curve $r=3+3\sin\theta.$\\

\input{polar.tikz}

\vspace{-4mm}

\clearpage\newpage
\probpts{13}{6}  Evaluate and simplify the indefinite integral \, $\ds \int e^{2x}\sin x\,dx$.
\vfill

\probpts{Extra Credit A (easier)}{2}  If you compute the millionth partial sum of the infinite series in problem {\large \textbf{7 (a)}}, about how many digits would you expect to be accurate?
\vfill


\clearpage\newpage
\probpts{Extra Credit B (harder)}{2}  Evaluate the integral in problem {\large \textbf{3 (c)}}.
\vfill


\clearpage\newpage
\centerline{\footnotesize \textsc{blank space}}
\vfill

\hrulefill

You may find the following \textbf{trigonometric formulas} useful.  Other formulas, not listed here, should be in your memory, or you can derive them from the ones here.
\begin{align*}
\sin(\alpha \pm \beta) &= \sin \alpha \cos \beta \pm \cos \alpha \sin \beta \qquad &
    \sin(ax) \sin(bx) &= \frac{1}{2} \cos((a-b)x) - \frac{1}{2} \cos((a+b)x) \\
\cos(\alpha \pm \beta) &= \cos \alpha \cos \beta \mp \sin \alpha \sin \beta \qquad &
    \sin(ax) \cos(bx) &= \frac{1}{2} \sin((a-b)x) + \frac{1}{2} \sin((a+b)x) \\
 && \cos(ax) \cos(bx) &= \frac{1}{2} \cos((a-b)x) + \frac{1}{2} \cos((a+b)x)
\end{align*}



\begin{comment}
%%%%%%%%%%%% commented out
\clearpage\newpage
\large\centerline{\textbf{Summary of Convergence Tests}}
\normalsize

\vspace{-2mm}
\begin{center}
\includegraphics[width=0.85\textwidth]{serieshandoutA.png}
\end{center}
\vfill

\clearpage\newpage
\begin{center}
\includegraphics[width=0.87\textwidth]{serieshandoutB.png}

\vspace{-2mm}
\includegraphics[width=0.87\textwidth]{serieshandoutC.png}
\end{center}
\vfill
\end{comment}

\end{document}